% takeoff_landing.tex -- Methodology: takeoff from the floor, lift and landing (draft 4 Oct 2026).
% Code: mpc_planner/creep_controller.py (arc creep, hand-over latch), planner_node.py (pretension,
% breakaway, ff_cap_release_s, lift ramp, LAND, TouchdownDetector), bringup/config/*.yaml.
% Evidence: 23 Sep rig non-lift (start_taut), 30 Sep rig review, R0817/R0818, R0834-R0839 (B2),
% R0807-R0816 (hand-over angle in the twin), RIG-1001-* (hand-over and rods on the rig).
% Condensed 12 Oct 2026 to the plan's half page; the detail is Appendix C "Floor start and landing in detail".
% Restructured 9 Oct (structure audit): the outcomes (the hand-over angle flown, fig:flight) moved to
% Section 4.1.

\subsubsection{Takeoff, lift and landing}
\label{sec:takeoff}

The planner assumes taut rods, but every flight starts and ends with the ring on the floor and the rods slack, where
the UAVs sit outside the cone the planner's references put them on; started from rest, the planner pulled them inwards
over the ring, the cause of the first failed lifts on the rig. On TAKEOFF each rod is therefore swept about its own
attachment point up to the planned elevation, 60$^\circ$ since 7 October, and the planner takes over once every rod
has held within 8$^\circ$ of it for 1\,s, measuring each rod's length at that hand-over. The planned pull then rises
over 3\,s with the ring still on the floor; the fraction reached when the ring breaks free is a cap on the trackers'
feedforward, released over 2\,s once the lift starts; and the ring's height target ramps up at 0.1\,m/s. On LAND the
target descends at 0.2\,m/s until the ring rests; in the twin each UAV then unwinds along its rod's circle (the rig's
tethers go slack and it lands straight down), the rod feedforward blends to hover over 1\,s, and a UAV back at its rest
idles armed. Touchdown is declared when the UAVs stop following the descending reference for 1\,s after at least
5\,cm of descent; in the twin this test is re-armed when the ring touches down (Section~\ref{sec:sim_landing}), and on
the rig after a review. Appendix~\ref{sec:app_floor} gives each step with its numbers and the failure it answers.
